% Computer Science A --- Advanced web fundamentals (HTML, CSS, JavaScript).
% Continuation of \texttt{cs-web/}; formatted from raw source in \texttt{cs-web-adv/temp.tex}.

\runningheader{Computer Science A}{Web: Advanced fundamentals}{Page \thepage\ of \numpages}

\begingradingrange{csawebadv}

\uplevel{\par\textbf{Part 1: HTML}\par}

\question[4]
TRUE or FALSE: An \texttt{id} attribute is a non-unique identifier that can be reused on multiple different elements within the same HTML page.
\begin{choices}
  \choice TRUE
  \CorrectChoice FALSE
\end{choices}
\begin{solution}
  Each \texttt{id} must be unique on a page. Reuse \texttt{class} for shared styling or grouping; reserve \texttt{id} for one-off targets (labels, anchors, \lstinline[style=jsinline]|document.getElementById|).
\begin{lstlisting}[style=htmlinline]
<button id="submit-btn" class="btn primary">Submit</button>
<button class="btn primary">Cancel</button>
\end{lstlisting}
  Two elements must not share \texttt{id="submit-btn"}, but both may share \texttt{class="btn primary"}.
\end{solution}

\question[4]
TRUE or FALSE: The DOM (Document Object Model) is a hierarchical tree structure that allows JavaScript to interact with and change HTML elements in real-time.
\begin{choices}
  \CorrectChoice TRUE
  \choice FALSE
\end{choices}
\begin{solution}
  The browser parses HTML into a tree rooted at \texttt{<html>}. JavaScript reads and mutates nodes to update the page without a full reload.
\begin{lstlisting}[style=htmlinline]
<h1 id="title">Hello</h1>
\end{lstlisting}
\begin{lstlisting}[style=js]
const title = document.getElementById('title');
title.textContent = 'Welcome back';
title.style.color = 'navy';
\end{lstlisting}
\end{solution}

\question[4]
TRUE or FALSE: HTML is a programming language used to perform complex logic and calculations within a web document.
\begin{choices}
  \choice TRUE
  \CorrectChoice FALSE
\end{choices}
\begin{solution}
  HTML is a markup language for structure and semantics. Logic, loops, and calculations belong in JavaScript (or another programming language).
\begin{lstlisting}[style=htmlinline]
<p>Items: <span id="count">0</span></p>
<button id="add">Add item</button>
\end{lstlisting}
\begin{lstlisting}[style=js]
let count = 0;
document.getElementById('add').onclick = () => {
  count += 1;
  document.getElementById('count').textContent = count;
};
\end{lstlisting}
  HTML declares what appears on the page; the script performs the calculation.
\end{solution}


\newpage


\question[4]
Which of the following is an example of a void tag? \emph{Select all that apply.}
\begin{checkboxes}
  \choice \texttt{<a>}
  \choice \texttt{<div>}
  \choice \texttt{<p>}
  \CorrectChoice \texttt{<img>}
\end{checkboxes}
\begin{solution}
  Void elements cannot enclose child nodes and have no closing tag. \texttt{<img>} is self-contained; \texttt{<a>}, \texttt{<div>}, and \texttt{<p>} wrap content and require closing tags.
\begin{lstlisting}[style=htmlinline]
<img src="avatar.png" alt="Profile photo">
<a href="/profile">View profile</a>
<p>Paragraph text goes here.</p>
\end{lstlisting}
  Other common void tags include \texttt{<br>}, \texttt{<input>}, and \texttt{<meta>}.
\end{solution}

\question[4]
TRUE or FALSE: To ensure accessibility, an \texttt{<input>} element should be connected to a \texttt{<label>} using the \texttt{for} attribute and a matching \texttt{id}.
\begin{choices}
  \CorrectChoice TRUE
  \choice FALSE
\end{choices}
\begin{solution}
  Pair \texttt{for} on the label with \texttt{id} on the input so screen readers announce the field and clicking the label text focuses the control.
\begin{lstlisting}[style=htmlinline]
<label for="user-email">Email Address:</label>
<input type="email" id="user-email" name="email">
\end{lstlisting}
\end{solution}

\question[4]
TRUE or FALSE: Block-level elements automatically start on a new line and take up the full width of their container.
\begin{choices}
  \CorrectChoice TRUE
  \choice FALSE
\end{choices}
\begin{solution}
  Block elements begin on a new line and stretch to the container width. Inline elements flow with surrounding text and only use the width their content needs.
\begin{lstlisting}[style=htmlinline]
<div>Block one (full width, new line)</div>
<div>Block two (full width, new line)</div>
<span>Inline A</span><span>Inline B</span>
\end{lstlisting}
  Common block elements: \texttt{<div>}, \texttt{<p>}, \texttt{<h1>}--\texttt{<h6>}, \texttt{<ul>}. Common inline elements: \texttt{<span>}, \texttt{<a>}, \texttt{<strong>}.
\end{solution}



\newpage



\question[4]
TRUE or FALSE: Semantic tags, such as \texttt{<nav>} and \texttt{<article>}, are used to provide meaning to the content for both search engines and assistive technologies.
\begin{choices}
  \CorrectChoice TRUE
  \choice FALSE
\end{choices}
\begin{solution}
  Semantic landmarks describe page regions explicitly. Crawlers and screen readers can identify navigation, main content, and articles instead of undifferentiated \texttt{<div>} wrappers.
\begin{lstlisting}[style=htmlinline]
<header><h1>Course site</h1></header>
<nav aria-label="Main">
  <a href="/">Home</a>
  <a href="/labs">Labs</a>
</nav>
<main>
  <article>
    <h2>Week 3 notes</h2>
    <p>Lesson content...</p>
  </article>
</main>
<footer>Contact info</footer>
\end{lstlisting}
\end{solution}

\question[4]
TRUE or FALSE: When using radio buttons, all buttons in the same group must share the same \texttt{name} attribute to ensure the user can only select one option.
\begin{choices}
  \CorrectChoice TRUE
  \choice FALSE
\end{choices}
\begin{solution}
  Radio inputs with the same \texttt{name} are mutually exclusive. Different \texttt{name} values are treated as separate groups.
\begin{lstlisting}[style=htmlinline]
<p>Preferred contact method:</p>
<input type="radio" id="email" name="contact-method" value="email">
<label for="email">Email</label>
<input type="radio" id="phone" name="contact-method" value="phone">
<label for="phone">Phone</label>
\end{lstlisting}
  Selecting \emph{Phone} automatically deselects \emph{Email} because both share \texttt{name="contact-method"}.
\end{solution}

\question[4]
Why is the \texttt{alt} attribute essential for an \texttt{<img>} tag? \emph{Select all that apply.}
\begin{checkboxes}
  \choice It defines the source path of the image file.
  \CorrectChoice It provides a text description for accessibility and SEO
  \choice It makes the image clickable as a hyperlink.
  \choice It allows the image to be resized automatically.
\end{checkboxes}
\begin{solution}
  \texttt{alt} gives alternate text for screen readers and when the image fails to load; search crawlers use it to index image content. The file path is \texttt{src}; hyperlinks wrap the image in \texttt{<a>}; sizing uses CSS or \texttt{width}/\texttt{height}.
\begin{lstlisting}[style=htmlinline]
<img src="sales-chart.png" alt="Bar chart of Q3 revenue by region">
<a href="/reports/q3">
  <img src="sales-chart.png" alt="Open the full Q3 report">
</a>
\end{lstlisting}
\end{solution}

\question[4]
TRUE or FALSE: The \texttt{<head>} section of an HTML document contains the visible content that users interact with on a webpage.
\begin{choices}
  \choice TRUE
  \CorrectChoice FALSE
\end{choices}
\begin{solution}
  \texttt{<head>} holds metadata; \texttt{<body>} holds visible, interactive content.
\begin{lstlisting}[style=htmlinline]
<!DOCTYPE html>
<html lang="en">
  <head>
    <meta charset="utf-8">
    <title>Lab 4</title>
    <link rel="stylesheet" href="styles.css">
  </head>
  <body>
    <h1>Visible heading users read and click under</h1>
    <button type="button">Start lab</button>
  </body>
</html>
\end{lstlisting}
\end{solution}

\newpage

\uplevel{\par\textbf{Part 2: CSS}\par}

\question[6]
What is a selector in web development? Describe the major kinds and how they are used.
\begin{solutionorbox}[5in]
\textbf{Definition:} In front-end development, a \textbf{selector} is a pattern that identifies DOM nodes so CSS can style them or JavaScript can manipulate them. The term also appears in UI frameworks (dropdown ``select'' components) and state libraries (Redux/NgRx \textbf{selectors}).

\textbf{Purpose:} Selectors decouple document structure from presentation and behavior: the same HTML can be restyled or scripted by changing the selector logic rather than rewriting markup.

\textbf{1. CSS selectors (styling):}
\begin{itemize}
  \item \textbf{Type/element} --- every instance of a tag: \lstinline[style=cssinline]|p|
  \item \textbf{Class} --- reusable groups: \lstinline[style=cssinline]|.btn-primary|
  \item \textbf{ID} --- one unique element per page: \lstinline[style=cssinline]|#main-header|
  \item \textbf{Attribute} --- match presence or value: \lstinline[style=cssinline]!input[type="email"]!
  \item \textbf{Pseudo-class / pseudo-element} --- state or fragment: \lstinline[style=cssinline]|button:hover|, \lstinline[style=cssinline]|p::first-line|
\end{itemize}

\begin{lstlisting}[style=css]
.btn-primary { background-color: steelblue; color: white; }
#main-header { font-size: 1.5rem; }
input[type="text"] { border: 1px solid gray; }
button:hover { background-color: darkblue; }
\end{lstlisting}

\textbf{2. DOM selectors (JavaScript):}
\begin{itemize}
  \item \lstinline[style=jsinline]!document.querySelector('.submit-btn')! --- first match for a CSS selector string.
  \item \lstinline[style=jsinline]!document.querySelectorAll('.alert')! --- static \lstinline[style=jsinline]|NodeList| of all matches (ideal for loops).
  \item Legacy: \lstinline[style=jsinline]|getElementById|, \lstinline[style=jsinline]|getElementsByClassName| --- faster in narrow cases but less expressive.
\end{itemize}

\begin{lstlisting}[style=js]
const submitBtn = document.querySelector('.submit-btn');
submitBtn.textContent = 'Click Me';

document.querySelectorAll('.alert-box')
  .forEach(alert => { alert.style.display = 'none'; });
\end{lstlisting}

\textbf{3. UI component selectors:} Native \lstinline[style=htmlinline]|<select>| dropdowns and custom combo-box/autocomplete widgets (for example react-select) let users pick one or many options inside forms.

\textbf{4. State-management selectors:} In Redux or NgRx, a selector is a pure function that derives a slice of global state (for example \lstinline[style=jsinline]|selectUserLoggedIn|) so components re-render only when that slice changes.
\end{solutionorbox}

\question[4]
When using Flexbox, which property is used to align items along the main axis? \emph{Select all that apply.}
\begin{checkboxes}
  \choice \texttt{align-items}
  \CorrectChoice \texttt{justify-content}
  \choice \texttt{flex-direction}
  \choice \texttt{align-content}
\end{checkboxes}
\begin{solution}
  \texttt{justify-content} distributes items along the main axis. \texttt{align-items} and \texttt{align-content} work on the cross axis; \texttt{flex-direction} sets the main axis direction but does not align items along it.
\begin{lstlisting}[style=css]
.toolbar {
  display: flex;
  flex-direction: row;      /* main axis is horizontal */
  justify-content: center;  /* align along main axis */
  align-items: center;      /* align along cross axis */
  gap: 1rem;
}
\end{lstlisting}
\end{solution}

\question[4]
TRUE or FALSE: The \texttt{rem} unit is relative to the font size of the parent element.
\begin{choices}
  \choice TRUE
  \CorrectChoice FALSE
\end{choices}
\begin{solution}
  \texttt{rem} is relative to the root (\texttt{<html>}) font size. \texttt{em} is relative to the parent (or the element's own font size for non-\texttt{font-size} properties).
\begin{lstlisting}[style=css]
html { font-size: 16px; }   /* browser default is often 16px */
.card { font-size: 1.25rem; }  /* 20px when root is 16px */
.card p { font-size: 1em; }    /* 20px: 1em of parent .card */
.sidebar { padding: 1.5rem; }  /* always 24px when root is 16px */
\end{lstlisting}
\end{solution}

\question[4]
TRUE or FALSE: A descendant selector targets only the immediate children of a parent element.
\begin{choices}
  \choice TRUE
  \CorrectChoice FALSE
\end{choices}
\begin{solution}
  A descendant selector (\texttt{.parent p}) matches any nested \texttt{p}. The child combinator (\texttt{.parent > p}) limits the match to direct children only.
\begin{lstlisting}[style=htmlinline]
<div class="parent">
  <p>Paragraph 1 (immediate child)</p>
  <div><p>Paragraph 2 (grandchild)</p></div>
</div>
\end{lstlisting}
\begin{lstlisting}[style=css]
.parent p  { color: navy; }  /* both paragraphs */
.parent > p { color: crimson; }  /* only Paragraph 1 */
\end{lstlisting}
\end{solution}

\question[4]
In CSS Grid, what does the \texttt{fr} unit represent? \emph{Select all that apply.}
\begin{checkboxes}
  \choice A fixed pixel value.
  \choice The font-size of the root element.
  \choice A percentage of the viewport width.
  \CorrectChoice A fraction of the available space in the grid container.
\end{checkboxes}
\begin{solution}
  \texttt{fr} divides leftover grid space after fixed tracks are placed. With \texttt{1fr 2fr 1fr}, the flexible area splits into four parts: side columns each get one part, the center gets two.
\begin{lstlisting}[style=css]
.dashboard {
  display: grid;
  grid-template-columns: 200px 1fr 2fr;
  gap: 1rem;
}
/* 200px sidebar, then remaining width split 1:2 */
\end{lstlisting}
\end{solution}

\question[4]
Which selector would you use to target only the very next sibling of an element? \emph{Select all that apply.}
\begin{checkboxes}
  \CorrectChoice \texttt{element + sibling}
  \choice \texttt{element \textasciitilde{} sibling}
  \choice \texttt{element > sibling}
  \choice \texttt{element sibling}
\end{checkboxes}
\begin{solution}
  The adjacent sibling combinator (\texttt{+}) matches the immediately following sibling with the same parent. General sibling (\texttt{{\textasciitilde}}) matches any later sibling; \texttt{>} selects children, not siblings.
\begin{lstlisting}[style=htmlinline]
<h2>Section title</h2>
<p>First paragraph after the heading.</p>
<p>Second paragraph (not adjacent to h2).</p>
\end{lstlisting}
\begin{lstlisting}[style=css]
h2 + p { margin-top: 0; font-weight: bold; }  /* only first p */
h2 ~ p { color: gray; }  /* both paragraphs */
\end{lstlisting}
\end{solution}

\question[4]
What are sibling selectors in CSS? How do the adjacent and general combinators differ?
\begin{solutionorbox}[3.75in]
\textbf{Definition:} \textbf{Sibling selectors} target elements that share the same parent and sit at the same nesting depth. CSS provides two combinators: \textbf{adjacent sibling} (\lstinline[style=cssinline]|+|) and \textbf{general sibling} (\lstinline[style=cssinline]!h1 ~ p!).

\textbf{Purpose:} Sibling selectors style relationships between peer elements without extra wrapper classes---common for intro paragraphs after headings, form field spacing, and CSS-only toggle/tab patterns.

\textbf{Adjacent sibling (\lstinline[style=cssinline]|+|):} Selects element $B$ only when $B$ is the \emph{immediate} next sibling after element $A$.

\begin{lstlisting}[style=html]
<h1>Main Title</h1>
<p>Styled: directly after the heading.</p>
<p>Not styled: second paragraph.</p>
\end{lstlisting}
\begin{lstlisting}[style=css]
h1 + p { color: navy; font-weight: bold; }  /* only first p */
\end{lstlisting}

\textbf{General sibling (\lstinline[style=cssinline]!h1 ~ p!):} Selects every $B$ that follows $A$ among shared-parent siblings, even if other elements appear between them.

\begin{lstlisting}[style=html]
<h1>Main Title</h1>
<p>Styled.</p>
<div>Unrelated container</div>
<p>Also styled.</p>
\end{lstlisting}
\begin{lstlisting}[style=css]
h1 ~ p { color: green; }  /* both paragraphs */
\end{lstlisting}

\textbf{Important rule:} Selectors only look \emph{forward} in document order. You cannot style a preceding sibling (for example an \lstinline[style=htmlinline]|<h1>| based on a later \lstinline[style=htmlinline]|<p>|).
\end{solutionorbox}

\question[4]
TRUE or FALSE: By default, adding padding to an element increases its total calculated width.
\begin{choices}
  \CorrectChoice TRUE
  \choice FALSE
\end{choices}
\begin{solution}
  With default \texttt{box-sizing: content-box}, \texttt{width} applies to content only; horizontal padding adds to the rendered width (\texttt{200px} content + \texttt{20px} padding each side = \texttt{240px} total).
\begin{lstlisting}[style=css]
.card {
  width: 200px;
  padding: 20px;           /* default content-box */
  box-sizing: content-box; /* browser default */
}
/* Many projects reset globally: */
*, *::before, *::after { box-sizing: border-box; }
\end{lstlisting}
  With \texttt{border-box}, padding stays inside the declared \texttt{200px}.
\end{solution}

\question[4]
TRUE or FALSE: In the CSS specificity hierarchy, an ID selector is more powerful than a class selector.
\begin{choices}
  \CorrectChoice TRUE
  \choice FALSE
\end{choices}
\begin{solution}
  ID selectors outrank classes in the specificity tuple \texttt{(inline, id, class, element)}. One \texttt{\#id} beats any number of concatenated class selectors.
\begin{lstlisting}[style=htmlinline]
<button id="save" class="btn primary">Save</button>
\end{lstlisting}
\begin{lstlisting}[style=css]
.btn.primary { background: gray; }  /* (0, 0, 2, 1) */
#save        { background: green; }  /* (0, 1, 0, 1) wins */
\end{lstlisting}
\end{solution}

\question[4]
Which unit is best suited for defining font sizes to ensure accessibility and scalability? \emph{Select all that apply.}
\begin{checkboxes}
  \choice \texttt{cm}
  \choice \texttt{px}
  \CorrectChoice \texttt{rem}
  \choice \texttt{pt}
\end{checkboxes}
\begin{solution}
  \texttt{rem} scales with the user's root font preference and browser zoom. Fixed \texttt{px}, \texttt{pt}, and \texttt{cm} do not respect text-size settings as reliably.
\begin{lstlisting}[style=css]
html { font-size: 100%; }     /* respects user browser setting */
body { font-size: 1rem; }     /* equals root size */
h1   { font-size: 2rem; }     /* doubles with root */
small { font-size: 0.875rem; }
\end{lstlisting}
\end{solution}

\question[4]
TRUE or FALSE: CSS transitions require a state change (like \texttt{:hover}) to trigger.
\begin{choices}
  \CorrectChoice TRUE
  \choice FALSE
\end{choices}
\begin{solution}
  A transition animates between two property values; something must change the style (\texttt{:hover}, \texttt{:focus}, or a class toggle). Use \texttt{@keyframes} for effects that run on load without a trigger.
\begin{lstlisting}[style=css]
.button {
  background: steelblue;
  transition: background 0.3s ease;
}
.button:hover,
.button:focus-visible {
  background: darkblue;  /* state change triggers transition */
}
\end{lstlisting}
\end{solution}

\question[4]
What is the primary advantage of a mobile-first strategy? \emph{Select all that apply.}
\begin{checkboxes}
  \choice It eliminates the need for media queries.
  \choice It makes websites look better on desktops.
  \CorrectChoice It results in cleaner code and better performance on mobile devices.
  \choice It uses only absolute units like \texttt{px}.
\end{checkboxes}
\begin{solution}
  Mobile-first starts with a minimal layout and adds complexity with \texttt{min-width} media queries, keeping default payloads light for constrained devices.
\begin{lstlisting}[style=css]
/* Base styles: single column for phones */
.grid { display: grid; grid-template-columns: 1fr; gap: 1rem; }

/* Enhance for wider screens */
@media (min-width: 768px) {
  .grid { grid-template-columns: 1fr 1fr; }
}
@media (min-width: 1024px) {
  .grid { grid-template-columns: 1fr 1fr 1fr; }
}
\end{lstlisting}
\end{solution}

\question[4]
What are media queries, and how do mobile-first and desktop-first strategies differ?
\begin{solutionorbox}[3.75in]
\textbf{Definition:} A \textbf{media query} is a CSS3 rule that applies styles only when device characteristics match a condition (screen width, orientation, resolution, or user preferences such as dark mode).

\textbf{Purpose:} Media queries are the foundation of \textbf{responsive web design}: one HTML document adapts layout and typography across phones, tablets, and desktops.

\textbf{Basic syntax:}
\begin{lstlisting}[style=css]
@media screen and (max-width: 768px) {
  body {
    background-color: #e0f2fe;
    font-size: 16px;
  }
}
\end{lstlisting}

\textbf{Common features inspected:}
\begin{itemize}
  \item Width and height: \lstinline[style=cssinline]|width|, \lstinline[style=cssinline]|min-width|, \lstinline[style=cssinline]|max-width|
  \item Orientation: \lstinline[style=cssinline]|portrait| vs.\ \lstinline[style=cssinline]|landscape|
  \item Resolution: high-density (Retina) displays
  \item User preferences: \lstinline[style=cssinline]|prefers-color-scheme: dark|, \lstinline[style=cssinline]|prefers-reduced-motion|
\end{itemize}

\textbf{Desktop-first (\lstinline[style=cssinline]|max-width|):} Start with large-screen styles, then add queries that tighten layout as the viewport shrinks.

\textbf{Mobile-first (\lstinline[style=cssinline]|min-width|) --- industry preference:} Ship minimal base styles for small screens (no query overhead), then progressively enhance with \lstinline[style=cssinline]|min-width| breakpoints as space grows. This keeps default payloads lighter on constrained devices.

\begin{lstlisting}[style=css]
/* Mobile-first base */
.container { width: 100%; }

@media screen and (min-width: 768px) {
  .container { width: 90%; }
}
@media screen and (min-width: 1024px) {
  .container { width: 80%; max-width: 1200px; }
}
\end{lstlisting}
\end{solutionorbox}

\newpage

\uplevel{\par\textbf{Part 3: JavaScript}\par}

\question[4]
TRUE or FALSE: Template literals require the use of \lstinline[style=jsinline]|\n| to create multi-line strings.
\begin{choices}
  \choice TRUE
  \CorrectChoice FALSE
\end{choices}
\begin{solution}
  Backtick template literals preserve literal line breaks. Only quoted strings need an explicit \lstinline[style=jsinline]|\n| escape.
\begin{lstlisting}[style=js]
const traditional = "Line one\nLine two";
const template = `Line one
Line two`;
console.log(template);
// Line one
// Line two
\end{lstlisting}
\end{solution}

\question[4]
TRUE or FALSE: The strict equality operator (\lstinline[style=jsinline]|===|) performs type coercion before comparing two values.
\begin{choices}
  \choice TRUE
  \CorrectChoice FALSE
\end{choices}
\begin{solution}
  \lstinline[style=jsinline]|===| compares value and type without coercion. \lstinline[style=jsinline]|==| may coerce types first.
\begin{lstlisting}[style=js]
console.log(5 === "5");  // false (number vs string)
console.log(5 === 5);    // true
console.log(5 == "5");   // true (string coerced to number)
console.log(null == undefined);  // true
console.log(null === undefined); // false
\end{lstlisting}
\end{solution}

\question[4]
TRUE or FALSE: The \lstinline[style=jsinline]|var| keyword is block-scoped, meaning it is only accessible within the \lstinline[style=jsinline]|{ }| where it was defined.
\begin{choices}
  \choice TRUE
  \CorrectChoice FALSE
\end{choices}
\begin{solution}
  \lstinline[style=jsinline]|var| is function-scoped and leaks out of blocks. \lstinline[style=jsinline]|let| and \lstinline[style=jsinline]|const| are block-scoped.
\begin{lstlisting}[style=js]
if (true) {
  var framework = "React";
  let library = "Zustand";
}
console.log(framework); // "React"
console.log(library);   // ReferenceError
\end{lstlisting}
\end{solution}

\question[4]
TRUE or FALSE: An arrow function inherits its \lstinline[style=jsinline]|this| value from its surrounding lexical scope.
\begin{choices}
  \CorrectChoice TRUE
  \choice FALSE
\end{choices}
\begin{solution}
  Arrow functions do not bind their own \lstinline[style=jsinline]|this|; they capture the enclosing scope's \lstinline[style=jsinline]|this|. Regular functions bind \lstinline[style=jsinline]|this| at call time.
\begin{lstlisting}[style=js]
const user = {
  name: "Jaehoon",
  greetRegular: function () {
    console.log(`Hello, ${this.name}`);
  },
  greetArrow: () => {
    console.log(`Hello, ${this.name}`);
  }
};
user.greetRegular(); // "Hello, Jaehoon"
user.greetArrow();   // "Hello, undefined" (inherits outer this)
\end{lstlisting}
\end{solution}

\question[4]
TRUE or FALSE: \lstinline[style=jsinline]|async| functions always return a \lstinline[style=jsinline]|Promise|.
\begin{choices}
  \CorrectChoice TRUE
  \choice FALSE
\end{choices}
\begin{solution}
  An \lstinline[style=jsinline]|async| function wraps return values in a resolved promise and thrown errors in a rejected promise, even when no explicit \lstinline[style=jsinline]|return| is written.
\begin{lstlisting}[style=js]
async function fetchStatus() {
  return "Success";
}
const result = fetchStatus();
console.log(result); // Promise { <fulfilled>: "Success" }

const value = await fetchStatus(); // "Success"
\end{lstlisting}
\end{solution}

\question[6]
Explain Promises, \lstinline[style=jsinline]|async|, and \lstinline[style=jsinline]|await|. How do they work together?
\begin{solutionorbox}[4.75in]
\textbf{Definition:} A \textbf{Promise} is an object representing the eventual result of an asynchronous operation (\emph{pending}, \emph{fulfilled}, or \emph{rejected}). The \lstinline[style=jsinline]|async| keyword marks a function that always returns a Promise; \lstinline[style=jsinline]|await| pauses execution inside that function until a Promise settles and yields its resolved value.

\textbf{Purpose:} Network calls, timers, and file I/O would block the main thread if written synchronously. Promises coordinate completion callbacks; \lstinline[style=jsinline]|async|/\lstinline[style=jsinline]|await| (ES2017) rewrites Promise chains into linear, readable code with standard \lstinline[style=jsinline]|try|/\lstinline[style=jsinline]|catch| error handling.

\textbf{Promise lifecycle with \lstinline[style=jsinline]|.then()| / \lstinline[style=jsinline]|.catch()|:}
\begin{lstlisting}[style=js]
function fetchUserData() {
  return new Promise((resolve, reject) => {
    const success = true;
    setTimeout(() => {
      if (success) {
        resolve({ id: 1, name: "Jaehoon" });
      } else {
        reject("Failed to fetch user");
      }
    }, 1000);
  });
}

fetchUserData()
  .then(data => console.log("User data received:", data))
  .catch(error => console.error("Error:", error));
\end{lstlisting}

\textbf{Modern \lstinline[style=jsinline]|async|/\lstinline[style=jsinline]|await| equivalent:}
\begin{lstlisting}[style=js]
async function displayUser() {
  try {
    console.log("Fetching user...");
    const data = await fetchUserData();
    console.log("User data received:", data);
  } catch (error) {
    console.error("Caught an error:", error);
  }
}

displayUser();
\end{lstlisting}

\textbf{Key contrasts:}
\begin{itemize}
  \item \textbf{Syntax:} \lstinline[style=jsinline]|.then()| chains callbacks; \lstinline[style=jsinline]|await| reads top-to-bottom like synchronous code.
  \item \textbf{Errors:} Promises use \lstinline[style=jsinline]|.catch()|; \lstinline[style=jsinline]|async| functions use \lstinline[style=jsinline]|try|/\lstinline[style=jsinline]|catch|.
  \item \textbf{Return behavior:} An \lstinline[style=jsinline]|async| function wraps non-Promise return values in an automatically resolved Promise; thrown errors become rejected Promises.
  \item \textbf{Under the hood:} \lstinline[style=jsinline]|async|/\lstinline[style=jsinline]|await| is syntactic sugar on top of Promises---it does not replace the Promise API.
\end{itemize}
\end{solutionorbox}

\question[4]
TRUE or FALSE: Array methods like \lstinline[style=jsinline]|.map()| and \lstinline[style=jsinline]|.filter()| mutate the original array they are called upon.
\begin{choices}
  \choice TRUE
  \CorrectChoice FALSE
\end{choices}
\begin{solution}
  \lstinline[style=jsinline]|.map()| and \lstinline[style=jsinline]|.filter()| return new arrays and leave the source unchanged. Mutating methods include \lstinline[style=jsinline]|.push()|, \lstinline[style=jsinline]|.pop()|, \lstinline[style=jsinline]|.splice()|, and \lstinline[style=jsinline]|.sort()|.
\begin{lstlisting}[style=js]
const numbers = [1, 2, 3, 4];
const evens = numbers.filter(n => n % 2 === 0);
const doubled = numbers.map(n => n * 2);

console.log(evens);    // [2, 4]
console.log(doubled);  // [2, 4, 6, 8]
console.log(numbers);  // [1, 2, 3, 4] unchanged
\end{lstlisting}
\end{solution}

\question[4]
TRUE or FALSE: You should commit the \lstinline[style=jsinline]|node_modules| folder to your Git repository to ensure other developers have the same dependencies.
\begin{choices}
  \choice TRUE
  \CorrectChoice FALSE
\end{choices}
\begin{solution}
  Track \lstinline[style=jsinline]|package.json| and the lockfile instead; add \lstinline[style=jsinline]|node_modules/| to \lstinline[style=bashinline]|.gitignore|. Teammates run \lstinline[style=bashinline]|npm install| locally.
\begin{lstlisting}[style=bash]
# .gitignore
node_modules/
\end{lstlisting}
\begin{lstlisting}[style=json]
{
  "name": "my-app",
  "dependencies": { "express": "^4.21.0" }
}
\end{lstlisting}
\begin{lstlisting}[style=bash]
npm install
\end{lstlisting}
\end{solution}

\question[4]
TRUE or FALSE: In JavaScript, a \lstinline[style=jsinline]|const| variable can have its reference modified if it holds an array or an object.
\begin{choices}
  \choice TRUE
  \CorrectChoice FALSE
\end{choices}
\begin{solution}
  \lstinline[style=jsinline]|const| locks the binding (reference) to one object or array. Internal properties or elements may still be mutated, but reassigning the variable throws a \texttt{TypeError}.
\begin{lstlisting}[style=js]
const user = { name: "Jaehoon" };
user.name = "Song";      // allowed: mutate contents
user = { name: "Alex" }; // TypeError: change reference

const ids = [1, 2];
ids.push(3);   // allowed: mutate array contents
ids = [4, 5];  // TypeError: change reference
\end{lstlisting}
\end{solution}

\question[4]
TRUE or FALSE: \lstinline[style=jsinline]|null| and \lstinline[style=jsinline]|undefined| are considered truthy values in JavaScript.
\begin{choices}
  \choice TRUE
  \CorrectChoice FALSE
\end{choices}
\begin{solution}
  Both are falsy. The full falsy set includes \lstinline[style=jsinline]|false|, \lstinline[style=jsinline]|0|, \lstinline[style=jsinline]|""|, \lstinline[style=jsinline]|null|, \lstinline[style=jsinline]|undefined|, and \lstinline[style=jsinline]|NaN|.
\begin{lstlisting}[style=js]
function isTruthy(value) {
  return value ? "truthy" : "falsy";
}
console.log(isTruthy(null));       // "falsy"
console.log(isTruthy(undefined));  // "falsy"
console.log(isTruthy([]));         // "truthy" (empty array)
console.log(isTruthy(""));         // "falsy"
\end{lstlisting}
\end{solution}

\question[4]
TRUE or FALSE: Event bubbling means that an event travels from the target element up through its ancestors to the root.
\begin{choices}
  \CorrectChoice TRUE
  \choice FALSE
\end{choices}
\begin{solution}
  After the target phase, the event propagates upward from the clicked node through parents to the document root. Call \lstinline[style=jsinline]|event.stopPropagation()| to halt bubbling.
\begin{lstlisting}[style=htmlinline]
<div id="panel">
  <button id="save">Save</button>
</div>
\end{lstlisting}
\begin{lstlisting}[style=js]
document.getElementById('panel').addEventListener('click', () => {
  console.log('panel heard click');
});
document.getElementById('save').addEventListener('click', (event) => {
  console.log('button clicked');
  // event.stopPropagation(); // uncomment to prevent panel log
});
// Clicking Save logs "button clicked" then "panel heard click"
\end{lstlisting}
\end{solution}

\endgradingrange{csawebadv}
